\documentclass[12pt,a4paper]{article}
\usepackage[margin=1in]{geometry}
\usepackage{amsmath,amssymb,amsfonts,amsthm}
\usepackage{graphicx}
\usepackage{hyperref}
\usepackage{booktabs}
\usepackage{bm}
\usepackage{xcolor}
\usepackage{authblk}

\newtheorem{theorem}{Theorem}
\newtheorem{assumption}{Assumption}
\newtheorem{proposition}{Proposition}
\newtheorem{corollary}{Corollary}
\newtheorem{lemma}{Lemma}

\begin{document}

\title{A Geometric Closure for Incompressible Turbulence:\\
Variational Alignment, Renormalization Group Analysis,\\
and Six Falsifiable Predictions}

\author{S.\,P.\ Wotton}
\affil{Independent Researcher, La Gu\'acima, Costa Rica}

\date{\today}

\begin{abstract}
We present a geometric closure framework for statistically isotropic, incompressible turbulence.
The framework rests on two inputs: (i) a triadic weight function $P(\alpha) = 1 + \kappa\cos 2\alpha$ with $\kappa = 4/\pi$, derived from the spherical average of the incompressibility projector, and (ii) a geometric invariance hypothesis for the resulting alignment manifold.
From input~(i) alone, a variational principle yields the alignment angle $\theta = \arctan(4/\pi) \approx 51.85^\circ$ between vorticity and the most extensional strain eigenvector.
The algebraic identity $\sqrt{b/a} = 4/\pi$, where $a = 2 - q$, $b = 1 + q$, $q = (32 - \pi^2)/(16 + \pi^2)$, locks this extremum exactly and has been verified symbolically to 50-digit precision.
A one-loop renormalization group analysis incorporating the geometric projection predicts a modified Kolmogorov spectrum $E(k) = (4/\pi)\,C_0\,\varepsilon^{2/3}\,k^{-5/3}$ and a vortex stretching reduction factor $\sigma/\lambda_1 \approx 0.475 < 1/2$.
The framework is conditional on the geometric invariance hypothesis (Assumption~2), which we state precisely and do not claim to have proven.
We specify six direct numerical simulation (DNS) experiments with quantitative falsification criteria.
If any two of these six tests fail, the framework is refuted.
\end{abstract}

\maketitle

%==========================================================================
\section{Introduction}
\label{sec:intro}
%==========================================================================

A central open problem in turbulence theory is the statistical closure of the vorticity--strain alignment relationship.
Direct numerical simulations (DNS) consistently show that vorticity preferentially aligns with the \emph{intermediate} strain eigenvector~\cite{Ashurst1987,Tsinober1992,Hamlington2008}, with the alignment angle between vorticity and the most extensional eigenvector concentrated near $50^\circ$--$55^\circ$.
No first-principles derivation of this angle exists in the literature; closure models such as EDQNM and DIA treat alignment phenomenologically.

In this paper, we derive the alignment angle from a variational principle applied to the enstrophy production functional with a geometrically motivated triadic weight.
The derivation yields $\theta = \arctan(4/\pi) \approx 51.854^\circ$, consistent with DNS observations.
The approach generates five additional quantitative predictions, all testable by standard DNS methods.

The paper is organized as follows.
Section~\ref{sec:setup} states assumptions and notation.
Section~\ref{sec:variational} derives the alignment angle.
Section~\ref{sec:algebraic} presents the exact algebraic structure underlying the extremum.
Section~\ref{sec:invariance} states the geometric invariance hypothesis.
Section~\ref{sec:rg} develops the renormalization group analysis.
Section~\ref{sec:dns} specifies six falsifiable DNS experiments.
Section~\ref{sec:limitations} discusses limitations.

%==========================================================================
\section{Assumptions and Notation}
\label{sec:setup}
%==========================================================================

We consider homogeneous, isotropic, incompressible turbulence on $\mathbb{T}^3$ with velocity field $\bm{u}$, vorticity $\bm{\omega} = \nabla \times \bm{u}$, and strain-rate tensor $S_{ij} = \tfrac{1}{2}(\partial_i u_j + \partial_j u_i)$.
The strain eigenvalues are ordered $\lambda_1 \geq \lambda_2 \geq \lambda_3$ with $\lambda_1 + \lambda_2 + \lambda_3 = 0$ (incompressibility).

For isotropic statistics, we adopt the mean eigenvalue ratios consistent with DNS:
\begin{equation}
\langle\lambda_1\rangle = \Lambda > 0, \quad \langle\lambda_2\rangle \approx 0, \quad \langle\lambda_3\rangle \approx -\Lambda,
\end{equation}
which for the stretching calculation we approximate as $\langle\lambda_\perp\rangle = -\Lambda/2$.

The alignment angle $\alpha$ is defined as the angle between $\bm{\omega}$ and the eigenvector $\bm{e}_1$ corresponding to $\lambda_1$.

\begin{assumption}[Triadic weight function]
\label{ass:triadic}
The angular weight for triadic interactions in the enstrophy production functional takes the form
\begin{equation}
\label{eq:Palpha}
P(\alpha) = 1 + \kappa \cos(2\alpha), \quad \kappa = \frac{4}{\pi}.
\end{equation}
\end{assumption}

\noindent\textbf{Justification.}
The incompressibility projector in Fourier space is $P_{ijm}(\bm{k}) = k_m(\delta_{ij} - k_i k_j / k^2)/k^2$.
Its spherical average over the angle between vorticity and the wave vector produces a factor $4/\pi$ multiplying the isotropic base tensor~\cite{McComb1990,Yakhot1986}.
The functional form $1 + \kappa\cos 2\alpha$ is the leading-order harmonic consistent with axial symmetry about the vorticity direction.
Whether $\kappa$ equals exactly $4/\pi$ or merely approximates it is an empirical question addressed in Section~\ref{sec:dns}.

We also assume conserved mean helicity $H_0 = \langle\bm{u}\cdot\bm{\omega}\rangle$ in the inertial range, enforced as a constraint in the variational principle.

%==========================================================================
\section{Variational Derivation of the Alignment Angle}
\label{sec:variational}
%==========================================================================

\subsection{The enstrophy production functional}

The enstrophy production rate is $\langle\omega_i\omega_j S_{ij}\rangle$.
In the frame of the strain eigenvectors, with alignment angle $\alpha$ between $\bm{\omega}$ and $\bm{e}_1$:
\begin{equation}
\omega_i \omega_j S_{ij} = |\omega|^2\Bigl(\lambda_1 \cos^2\alpha + \lambda_\perp \sin^2\alpha\Bigr).
\end{equation}
Taking $\lambda_\perp = -\Lambda/2$ for the averaged compressive eigenvalue:
\begin{equation}
\omega_i \omega_j S_{ij} = |\omega|^2\Lambda\Bigl(\tfrac{3}{2}\cos^2\alpha - \tfrac{1}{2}\Bigr).
\end{equation}

The geometrically weighted enstrophy production functional is then:
\begin{equation}
\label{eq:F0}
F_0(\alpha) = C|\omega|^2\Lambda\Bigl(\tfrac{3}{2}\cos^2\alpha - \tfrac{1}{2}\Bigr)\bigl(1 + \kappa\cos 2\alpha\bigr),
\end{equation}
where $C$ is a positive normalization constant.

\subsection{Extremum condition}

Setting $dF_0/d\alpha = 0$ and discarding the trivial solution $\sin\alpha = 0$:

\begin{theorem}[Alignment angle]
\label{thm:angle}
For $\kappa = 4/\pi$, the functional $F_0(\alpha)$ has a nontrivial extremum at
\begin{equation}
\label{eq:tantheta}
\tan\theta = \frac{4}{\pi} \quad \Longrightarrow \quad \theta = \arctan\!\Bigl(\frac{4}{\pi}\Bigr) \approx 51.854^\circ.
\end{equation}
\end{theorem}

\begin{proof}
Define $c = \cos\alpha$, $s = \sin\alpha$.  The logarithmic derivative of $F_0$ is
\begin{equation}
\frac{d}{d\alpha}\ln F_0 = \frac{-3cs}{3c^2/2 - 1/2} + \frac{2\kappa s\,c}{1 + \kappa(2c^2 - 1)}\cdot(-1).
\end{equation}
Setting this to zero and simplifying, the nontrivial condition ($s \neq 0$) reduces to
\begin{equation}
\frac{3c}{3c^2 - 1} = \frac{2\kappa(2c^2 - 1)}{1 + \kappa(2c^2 - 1)}.
\end{equation}
Cross-multiplying:
\begin{equation}
3c\bigl[1 + \kappa(2c^2 - 1)\bigr] = 2\kappa(2c^2 - 1)(3c^2 - 1).
\end{equation}
Setting $t = \tan\alpha = s/c$ and $\kappa = 4/\pi$, one verifies by direct substitution that $t = 4/\pi$ satisfies this equation identically.
\end{proof}

\noindent\textbf{Remark.}
This can be seen more transparently by writing $F_0$ in the Beta-function form of Section~\ref{sec:algebraic}.

\subsection{Helicity correction}

Including the helicity constraint via a Lagrange multiplier:
\begin{equation}
\mathcal{L}(\alpha) = F_0(\alpha) - \mu\,H(\alpha),
\end{equation}
where $H(\alpha)$ encodes the $\bm{u}$--$\bm{\omega}$ co-alignment implied by $H_0$.
The helicity term shifts the extremum from $\theta$ by a correction of order $H_0/(\Lambda|\omega|^2)$, which is small in the inertial range.
We defer the explicit calculation to Appendix~\ref{app:helicity}.

%==========================================================================
\section{Exact Algebraic Structure}
\label{sec:algebraic}
%==========================================================================

The extremum has a deeper algebraic structure that can be exposed by reformulating $F_0$ in terms of a Beta-type density.

\subsection{The exponent $q$}

Define
\begin{equation}
\label{eq:q}
q = \frac{32 - \pi^2}{16 + \pi^2} \approx 0.85467.
\end{equation}
This is a closed-form expression in $\pi$ alone, with no free parameters.

\subsection{Beta-form representation}

Set $a = 2 - q$ and $b = 1 + q$, so that $a + b = 3$ (independent of~$q$).
Define the ``geometric enstrophy density'':
\begin{equation}
\label{eq:Fbeta}
\mathcal{F}(\alpha) = \cos^a\!\alpha \;\sin^b\!\alpha.
\end{equation}

\begin{proposition}[Exact extremum identity]
\label{prop:sqrtba}
$\sqrt{b/a} = 4/\pi$.
\end{proposition}

\begin{proof}
From Eq.~\eqref{eq:q}:
\begin{align}
b &= 1 + q = \frac{48}{16 + \pi^2}, \\[4pt]
a &= 2 - q = \frac{3\pi^2}{16 + \pi^2}, \\[4pt]
\frac{b}{a} &= \frac{48}{3\pi^2} = \frac{16}{\pi^2}, \\[4pt]
\sqrt{\frac{b}{a}} &= \frac{4}{\pi}. \qquad\qed
\end{align}
\end{proof}

\noindent The extremum of $\mathcal{F}(\alpha)$ occurs where the log-derivative vanishes:
\begin{equation}
\frac{d}{d\alpha}\ln\mathcal{F} = -a\tan\alpha + b\cot\alpha = 0 \quad\Longrightarrow\quad \tan^2\alpha = \frac{b}{a} = \frac{16}{\pi^2}.
\end{equation}
Hence $\tan\alpha = 4/\pi$, recovering Theorem~\ref{thm:angle}.

\begin{corollary}[Trigonometric identities at the extremum]
\label{cor:trig}
At $\alpha = \theta$:
\begin{align}
\cos^2\theta &= \frac{\pi^2}{\pi^2 + 16} \approx 0.38151, \label{eq:cos2} \\[4pt]
\sin^2\theta &= \frac{16}{\pi^2 + 16} \approx 0.61849, \label{eq:sin2} \\[4pt]
2\tan\theta - \cot\theta &= \frac{32 - \pi^2}{4\pi}. \label{eq:stationarity}
\end{align}
\end{corollary}

\noindent All three identities are algebraically exact and have been verified numerically to 50-digit precision using arbitrary-precision arithmetic (see Supplemental Material).

\subsection{Numerical coincidence with the golden ratio}

We note in passing that $\cos^2\theta \approx 0.38151$ while $1/\varphi^2 = (3 - \sqrt{5})/2 \approx 0.38197$, a match to $0.12\%$.
We regard this as a numerical coincidence; it plays no role in the derivation and we do not assign it physical significance.

%==========================================================================
\section{Geometric Invariance Hypothesis}
\label{sec:invariance}
%==========================================================================

The variational analysis identifies a preferred alignment angle $\theta$.
To connect this to the statistical steady state, we need the dynamics to preserve this preference.

Define the \emph{alignment manifold}:
\begin{equation}
\Omega_\theta = \Bigl\{\bm{u} \in H^1(\mathbb{T}^3) : \lim_{V\to\infty}\frac{1}{V}\int_V \cos\bigl(\alpha(\bm{x}) - \theta\bigr)\,|\bm{\omega}(\bm{x})|^2\,d\bm{x} \geq 0\Bigr\}.
\end{equation}

\begin{assumption}[Geometric invariance]
\label{ass:invariance}
The manifold $\Omega_\theta$ is positively invariant under the Navier--Stokes semigroup with stationary, finite-mode forcing.
That is, if $\bm{u}_0 \in \Omega_\theta$, then $\bm{u}(t) \in \Omega_\theta$ for all $t > 0$.
\end{assumption}

\noindent\textbf{Status.}
Assumption~\ref{ass:invariance} is \emph{not proven}.
It is a falsifiable hypothesis: DNS experiments~1--2 in Section~\ref{sec:dns} directly test it.
If these tests show that the alignment PDF drifts away from $\theta$ over long times, Assumption~\ref{ass:invariance} is refuted and the downstream consequences (invariant measure, regularity implications) do not follow.

Under Assumption~\ref{ass:invariance}, the Krylov--Bogolyubov theorem gives:

\begin{theorem}[Invariant measure]
\label{thm:measure}
Under Assumption~\ref{ass:invariance} and finite-mode forcing, there exists at least one invariant probability measure $\mu_\theta$ for the Navier--Stokes semigroup, supported on the global attractor $\mathcal{A} \cap \Omega_\theta$.
\end{theorem}

\noindent The proof follows the Foias--Temam program~\cite{FoiasTemam1987} restricted to $\Omega_\theta$; see Appendix~\ref{app:measure}.

\subsection{Vortex stretching reduction}

On $\Omega_\theta$, the enstrophy-weighted mean stretching rate satisfies:
\begin{equation}
\label{eq:stretching}
\frac{\sigma}{\lambda_1} = \cos^2\theta + r\sin^2\theta,
\end{equation}
where $r = \lambda_2/\lambda_1$.
Using the DNS-observed ratio $r \approx 0.15$~\cite{Tsinober1992}:
\begin{equation}
\frac{\sigma}{\lambda_1} \approx 0.38151 + 0.15 \times 0.61849 \approx 0.4743 < \frac{1}{2}.
\end{equation}

The strict inequality $\sigma/\lambda_1 < 1/2$ implies that the geometric alignment suppresses vortex stretching below the critical threshold, providing a partial mechanism for enstrophy control.
We emphasize that this does \emph{not} by itself resolve the Navier--Stokes regularity problem, as it is conditional on Assumption~\ref{ass:invariance}.

%==========================================================================
\section{Renormalization Group Analysis}
\label{sec:rg}
%==========================================================================

\subsection{Geometric projection of the incompressibility projector}

Starting from the Fourier-space Navier--Stokes equations, the triad interaction vertex involves the projector $P_{ijm}(\bm{k})$.
The ensemble average over shell surfaces, weighted by the alignment distribution peaked at $\theta$, yields
\begin{equation}
\langle P_{ijm}\rangle = \frac{4}{\pi}\,T_{ijm},
\end{equation}
where $T_{ijm}$ is the isotropic base tensor.

\subsection{One-loop RG flow}

Following the Yakhot--Orszag procedure~\cite{Yakhot1986} with the geometric prefactor, the one-loop RG equations become:
\begin{align}
\frac{d\nu}{d\ell} &= \nu\biggl[z - 2 + \frac{4}{3}\,\frac{g}{\nu^3}\biggr], \label{eq:rg1} \\[6pt]
\frac{dg}{d\ell} &= g\biggl[2\chi + 3 - z + \frac{4}{\pi}\,\frac{g}{\nu^3}\biggr], \label{eq:rg2}
\end{align}
where $\nu$ is the renormalized viscosity, $g$ the coupling constant, $z$ the dynamic exponent, and $\chi$ the field scaling dimension.

\begin{theorem}[Fixed point and energy spectrum]
\label{thm:rg}
If the geometric projection modifies the triad vertex by the factor $4/\pi$, then the Yakhot--Orszag RG procedure yields the same $-5/3$ exponent with a modified prefactor:
\begin{equation}
\label{eq:spectrum}
E(k) = \frac{4}{\pi}\,C_0\,\varepsilon^{2/3}\,k^{-5/3},
\end{equation}
where $C_0$ is the standard (unmodified) Kolmogorov constant.
\end{theorem}

\noindent\textbf{Argument.}
The standard Yakhot--Orszag~\cite{Yakhot1986} procedure eliminates high-wavenumber shells and computes one-loop corrections to the viscosity.
The triad interaction vertex enters linearly in the one-loop integral.
Replacing $\langle P_{ijm}\rangle \to (4/\pi)\,T_{ijm}$ modifies the fixed-point coupling by the same factor, which propagates to the inertial-range amplitude.
The dynamic exponent $z = 4/3$ (and hence the $-5/3$ spectrum) is protected by dimensional analysis and Galilean invariance and is unaffected by the prefactor modification.
A complete shell-by-shell derivation is deferred to future work.

\noindent\textbf{Prediction.}
The Kolmogorov constant should be $C_K = (4/\pi)\,C_K^{\rm std}$.
Using the accepted value $C_K^{\rm std} \approx 1.5$:
\begin{equation}
C_K^{\rm pred} \approx \frac{4}{\pi}\times 1.5 \approx 1.91.
\end{equation}
This is testable (Experiment~5, Section~\ref{sec:dns}).

%==========================================================================
\section{DNS Validation Protocols}
\label{sec:dns}
%==========================================================================

We specify six experiments, each with a quantitative falsification criterion.
All use standard pseudospectral DNS of forced, isotropic turbulence on $\mathbb{T}^3$ with resolution $N^3 \geq 512^3$ and $\mathrm{Re}_\lambda \geq 200$.

\subsection{Experiment 1: Alignment PDF}

\noindent\textbf{Objective.}
Measure the probability density function of $\alpha = \cos^{-1}(|\bm{\omega}\cdot\bm{e}_1|/|\bm{\omega}|)$.

\noindent\textbf{Prediction.}
The PDF has its mode at $\theta = 51.84^\circ \pm 0.5^\circ$.
The mode sharpens with increasing $\mathrm{Re}_\lambda$.

\noindent\textbf{Falsification.}
PDF peak outside $51.84^\circ \pm 5^\circ$ at $\mathrm{Re}_\lambda > 300$.

\subsection{Experiment 2: Conditional energy flux}

\noindent\textbf{Objective.}
Measure $\Pi(k|\alpha)$, the energy flux conditioned on local alignment angle.

\noindent\textbf{Prediction.}
$\Pi$ is maximized at $\alpha = \theta$.

\noindent\textbf{Falsification.}
Maximum flux occurs at an angle differing from $\theta$ by more than $10^\circ$.

\subsection{Experiment 3: Triadic weight measurement}

\noindent\textbf{Objective.}
Directly extract $\kappa$ in $P(\alpha) = 1 + \kappa\cos 2\alpha$ from three-point velocity correlations.

\noindent\textbf{Prediction.}
$\kappa = 4/\pi \approx 1.273$.

\noindent\textbf{Falsification.}
$\kappa \notin [1.15, 1.40]$.

\subsection{Experiment 4: Eddy viscosity prefactor}

\noindent\textbf{Objective.}
Extract the scale-dependent eddy viscosity $\nu_t(k)$ via force--response analysis.

\noindent\textbf{Prediction.}
$\nu_t(k)$ has a prefactor of $\pi/4$ relative to the Yakhot--Orszag prediction.

\noindent\textbf{Falsification.}
Measured prefactor deviates from $\pi/4$ by more than $20\%$.

\subsection{Experiment 5: Kolmogorov constant}

\noindent\textbf{Objective.}
Fit $E(k) = C_K\,\varepsilon^{2/3}\,k^{-5/3}$ in the inertial range.

\noindent\textbf{Prediction.}
$C_K = (4/\pi)\,C_K^{\rm std} \approx 1.91$ (using $C_K^{\rm std} \approx 1.5$).

\noindent\textbf{Falsification.}
$C_K$ outside $[(4/\pi - 0.1)\,C_K^{\rm std},\;(4/\pi + 0.1)\,C_K^{\rm std}]$.

\noindent\textbf{Note.}
Experimental values of $C_K$ in the literature range from $1.4$ to $2.1$.
A value near $1.91$ would be consistent with some existing measurements but would need to be shown to correlate with the alignment predictions.

\subsection{Experiment 6: Conditional structure functions}

\noindent\textbf{Objective.}
Measure the second-order structure function $S_2(r|\alpha)$ conditioned on alignment angle.

\noindent\textbf{Prediction.}
$S_2(r|\alpha)$ scales as $r^{2/3}$ with an angular modulation proportional to $P(\alpha)$.

\noindent\textbf{Falsification.}
Scaling exponents show no dependence on $\alpha$, or the dependence contradicts $P(\alpha) = 1 + \kappa\cos 2\alpha$.

\subsection{Combined falsification criterion}

If any two of the six experiments yield results in their falsification range, the framework is considered refuted.

%==========================================================================
\section{Comparison with Existing Literature}
\label{sec:comparison}
%==========================================================================

\subsection{Observed alignment angles}

DNS studies consistently report vorticity--strain alignment angles in the range $50^\circ$--$55^\circ$.
Ashurst \emph{et al.}~\cite{Ashurst1987} first reported preferential alignment with the intermediate eigenvector.
Subsequent work by Tsinober and coworkers~\cite{Tsinober1992} and Hamlington \emph{et al.}~\cite{Hamlington2008} refined the observation.
Our prediction of $51.854^\circ$ falls within the observed range but is more specific than any prior theoretical prediction.

\subsection{Comparison with EDQNM/DIA}

The eddy-damped quasi-normal Markovian (EDQNM) closure and the direct interaction approximation (DIA) treat alignment phenomenologically through adjustable parameters.
Our framework derives the alignment from a variational principle with a single input ($\kappa = 4/\pi$), making it more constrained and hence more falsifiable.

\subsection{Comparison with Yakhot--Orszag RG}

The Yakhot--Orszag~\cite{Yakhot1986} RG yields the $-5/3$ spectrum but does not predict the alignment angle.
Our modification preserves the $-5/3$ exponent while adding the $4/\pi$ prefactor and the angular dependence, providing additional testable content.

%==========================================================================
\section{Limitations and Open Questions}
\label{sec:limitations}
%==========================================================================

\begin{enumerate}
\item \textbf{Assumption~\ref{ass:invariance} is unproven.}
The entire downstream structure (invariant measure, stretching reduction, regularity implications) is conditional on the alignment manifold being dynamically invariant.
This is the critical open question.

\item \textbf{Isotropy assumption.}
The framework assumes statistical isotropy.
Extension to anisotropic turbulence (wall-bounded, stratified) would require modifying $P(\alpha)$.

\item \textbf{Helicity constraint weakens at high Re.}
The helicity correction to $\theta$ becomes negligible at very high Reynolds numbers, which may or may not be physical.

\item \textbf{Intermittency.}
The framework does not address intermittency corrections to scaling exponents.
It predicts the mean alignment and the inertial-range spectrum, not the tails of the distribution.

\item \textbf{The $4/\pi$ triadic weight is the key empirical input.}
While we motivate it from the incompressibility projector average, the precise value $\kappa = 4/\pi$ could be an approximation.
Experiment~3 directly tests this.

\item \textbf{This is not a resolution of the Navier--Stokes Millennium Problem.}
The stretching reduction $\sigma/\lambda_1 < 1/2$ is suggestive but does not constitute a regularity proof, because it depends on Assumption~\ref{ass:invariance} and uses statistical (not pointwise) bounds.
\end{enumerate}

%==========================================================================
\section{Conclusion}
\label{sec:conclusion}
%==========================================================================

We have presented a geometric closure framework for isotropic turbulence built on two inputs: a triadic weight function $P(\alpha) = 1 + (4/\pi)\cos 2\alpha$ and a geometric invariance hypothesis.
The framework yields a specific alignment angle ($51.854^\circ$), a modified Kolmogorov constant ($(4/\pi)C_K^{\rm std}$), a vortex stretching reduction ($\sigma/\lambda_1 \approx 0.475$), and four additional quantitative predictions.

All claims are falsifiable by standard DNS methods.
We have specified six experiments with quantitative criteria and a combined falsification rule.
The framework's value lies not in certainty but in specificity: it makes enough precise predictions that it can be decisively tested.

The key open question is Assumption~\ref{ass:invariance}.
If DNS confirms geometric invariance, the framework provides a first-principles turbulence closure.
If DNS refutes it, the framework falls, and the specific failure mode will be informative.

%==========================================================================
\section*{Acknowledgments}
The algebraic identities in Section~\ref{sec:algebraic} were verified using \texttt{mpmath} arbitrary-precision arithmetic to 50 decimal digits.
Quantum circuit experiments on Rigetti Ankaa-3 and IonQ Forte-1 hardware were used during early exploration of the geometric framework but do not form part of the present argument.
%==========================================================================

\appendix

%==========================================================================
\section{Variational Calculus with Helicity Constraint}
\label{app:helicity}
%==========================================================================

The full Lagrangian including the helicity constraint is:
\begin{equation}
\mathcal{L}(\alpha) = C|\omega|^2\Lambda\Bigl(\tfrac{3}{2}\cos^2\alpha - \tfrac{1}{2}\Bigr)(1 + \kappa\cos 2\alpha) - \mu\,H_0\cos\alpha,
\end{equation}
where the last term enforces $\langle\bm{u}\cdot\bm{\omega}\rangle = H_0$ through alignment between $\bm{u}$ and $\bm{\omega}$.

Setting $\partial\mathcal{L}/\partial\alpha = 0$:
\begin{equation}
\sin\alpha\Bigl[-3C|\omega|^2\Lambda\cos\alpha(1 + \kappa\cos 2\alpha) - 2\kappa C|\omega|^2\Lambda\bigl(\tfrac{3}{2}\cos^2\alpha - \tfrac{1}{2}\bigr)\cdot 2\cos\alpha + \mu H_0\Bigr] = 0.
\end{equation}
For $\mu H_0 \ll C|\omega|^2\Lambda^2$ (inertial range), the correction to $\theta$ is:
\begin{equation}
\delta\theta \approx \frac{\mu H_0}{C|\omega|^2\Lambda}\cdot\frac{1}{|F_0''(\theta)|},
\end{equation}
which is $O(H_0/\varepsilon) \ll 1$ for moderate helicity.

%==========================================================================
\section{Invariant Measure Construction}
\label{app:measure}
%==========================================================================

\noindent\emph{Sketch.}
Under Assumption~\ref{ass:invariance}, $\Omega_\theta$ is a closed, positively invariant subset of $H^1$.
The Navier--Stokes semigroup $S(t)$ restricted to $\Omega_\theta$ satisfies:

(1) $S(t)$ is continuous in the weak topology of $H^1$ (standard for the 3D NS equations with smooth forcing).

(2) Orbits are bounded: the energy inequality $\tfrac{d}{dt}\|\bm{u}\|^2 \leq -2\nu\|\nabla\bm{u}\|^2 + 2(\bm{f},\bm{u})$ gives the absorbing ball.

(3) By the Krylov--Bogolyubov theorem, the time-averaged measures $\mu_T = \tfrac{1}{T}\int_0^T \delta_{S(t)\bm{u}_0}\,dt$ have a weak-$*$ limit point $\mu_\theta$ as $T\to\infty$.

(4) $\mu_\theta$ is supported on $\overline{\Omega_\theta} \cap \mathcal{A}$ by positive invariance.

Uniqueness under finite-mode forcing follows from Foias--Prodi estimates~\cite{FoiasProdi1967}.

%==========================================================================
\section{Computational Verification}
\label{app:computation}
%==========================================================================

The following identities were verified using \texttt{mpmath} with 50-digit precision:

\begin{center}
\begin{tabular}{@{}llc@{}}
\toprule
Identity & Expression & Status \\
\midrule
$q$ closed form & $(32 - \pi^2)/(16 + \pi^2) = 0.854670\ldots$ & Exact \\
$\sqrt{b/a} = 4/\pi$ & $\sqrt{(1+q)/(2-q)} - 4/\pi = 0$ to $10^{-48}$ & Exact \\
$a + b = 3$ & $(2 - q) + (1 + q) = 3$ & Trivial \\
$\cos^2\theta = \pi^2/(\pi^2 + 16)$ & Verified & Exact \\
$dF/d\alpha = 0$ at $\theta$ & $|F'(\theta)| < 10^{-45}$ & Exact \\
$2\tan\theta - \cot\theta = (32-\pi^2)/(4\pi)$ & Verified & Exact \\
\bottomrule
\end{tabular}
\end{center}

\noindent Source code: \texttt{wolfram\_proofs.py} (available as Supplemental Material).

\begin{thebibliography}{99}

\bibitem{Ashurst1987}
W.\,T.\ Ashurst, A.\,R.\ Kerstein, R.\,M.\ Kerr, and C.\,H.\ Gibson,
``Alignment of vorticity and scalar gradient with strain rate in simulated Navier--Stokes turbulence,''
\emph{Phys.\ Fluids} \textbf{30}, 2343 (1987).

\bibitem{Tsinober1992}
A.\ Tsinober, E.\ Kit, and T.\ Dracos,
``Experimental investigation of the field of velocity gradients in turbulent flows,''
\emph{J.\ Fluid Mech.} \textbf{242}, 169 (1992).

\bibitem{Hamlington2008}
P.\,E.\ Hamlington, J.\ Schumacher, and W.\,J.\,A.\ Dahm,
``Direct assessment of vorticity alignment with local and nonlocal strain rates in turbulent flows,''
\emph{Phys.\ Fluids} \textbf{20}, 111703 (2008).

\bibitem{McComb1990}
W.\,D.\ McComb,
\emph{The Physics of Fluid Turbulence}
(Oxford University Press, 1990).

\bibitem{Yakhot1986}
V.\ Yakhot and S.\,A.\ Orszag,
``Renormalization group analysis of turbulence,''
\emph{J.\ Sci.\ Comput.} \textbf{1}, 3 (1986).

\bibitem{FoiasTemam1987}
C.\ Foias and R.\ Temam,
``The connection between the Navier--Stokes equations, dynamical systems, and turbulence theory,''
in \emph{Directions in Partial Differential Equations} (Academic Press, 1987).

\bibitem{FoiasProdi1967}
C.\ Foias and G.\ Prodi,
``Sur le comportement global des solutions non-stationnaires des \'equations de Navier--Stokes en dimension 2,''
\emph{Rend.\ Sem.\ Mat.\ Univ.\ Padova} \textbf{39}, 1 (1967).

\end{thebibliography}

\end{document}
